\documentclass[a4paper,10pt]{scrartcl}

\usepackage[ngerman]{babel}
\usepackage[utf8]{inputenc}
\usepackage[T1]{fontenc}
\usepackage[babel,german=quotes]{csquotes}
\usepackage{hyperref}
\usepackage{geometry}
\usepackage{graphicx}
\usepackage{enumitem}
\usepackage{multicol}

\setlength\parindent{0pt}
\setitemize{itemsep=0pt}


\geometry{
	a4paper,
	left=20mm,
	top=20mm,
	bottom=15mm,
	footskip=5mm
}



\title{Entwicklerhandbuch für das Thema \glqq\textit{Klick-Tool}\grqq}
\subtitle{Betreuer: \textit{Prof. Dr. Andreas Nüchter}}
\author{\textit{Tom Lerzer, Fabia Frank, Luis Ángel Álvarez Gil}}
\date{\textit{SS2026}}


\begin{document}

\maketitle

\tableofcontents

\newpage

\section{Beschreibung}

Die in der Programmiersprache C++ entwickelte Software \textit{new\_click\_tool} stellt ein Subprojekt im Projekt \textit{3DTK} der Universität Würzburg dar. Das Tool dient der Verarbeitung von 2D-Bildern, indem Punkte in einzelnen Bildern sowie Punktkorrespondenzen zwischen zwei Bildern markiert werden können. Die festgelegten Punkte werden in Koordinatenform als Textdatei ausgegeben. Hierzu wurde eine Zoom-Funktion implementiert um das Anwählen in Subpixelgenauigkeit von mindestens 0,1 Pixeln zu gewährleisten. Im Zusammenspiel mit anderen 3DTK-Subfunktionen können auch 3D-Formate eingegeben und intern zur weiteren Verarbeitung umgewandelt werden. 

\newpage

\section{Frontend / Oberfläche}

\subsection{Übersicht}

Das Frontend des \textit{new\_click\_tool} wird durch Nutzung der Grafikbibliothek \textit{Dear ImGui} implementiert. Es besteht aus einem Hauptfenster, in dem ein oder zwei Bilder angezeigt, vergrößert und bearbeitet werden können sowie einem Fenster \textit{Control \& Setup}, in dem Funktionalitäten an- und abgewählt werden können. Das \textit{Control \& Setup}-Fenster ermöglicht weiterhin das Laden von Bildern, die Konvertierung von 3D-Formaten in 2D-Bilder sowie das Speichern gewählter Punkte bzw. Punktkorrespondenzen.


\subsection{Datei-/Ordnerstruktur}
\begin{verbatim}
	3DTK/
			|--src/slam6d/fbr/new_click_tool/
			|		|--docs/
			|		|		|--Benutzerhandbuch.pdf
			|		|		|--Entwicklerhandbuch.pdf
			|		|		|--Pflichtenheft.pdf
			|		|--MVP /
			|		|		|--assets/
			|		|		|		|--mein_Bild.jpg
			|		|		|		|--...
			|		|		|--CMakeLists.txt
			|		|		|--new_click_tool_app.cpp
			|		|		|--new_click_tool_app_two.cpp
			|		|		|--new_click_tool_main.cpp
			|		|		|--...
			|--include/slam6d/fbr
			|		|--...
			|		|--new_click_tool_app.h
			|		|--...
			|--bin/
			|		|--...
			|		|--new_click_tool
			|		|--...
			
\end{verbatim}

\begin{itemize}
	\item Im Ordner \textbf{src/slam6d/fbr/new\_click\_tool/docs} befindet sich die im Rahmen des Projektes erstellte Dokumentation
	\item Der Ordner \textbf{new\_click\_tool/MVP} enthält, neben Beispielbildern in \textbf{assets}, den Quellcode des Tools. Dieser besteht u.a. aus:
	\begin{itemize}
		\item new\_click\_tool\_app.cpp: Wesentliche Bestandteile der Benutzeroberfläche sowie Zoom-Funktion und Bilddarstellung
		\item new\_click\_tool\_app\_two.cpp: Erweiterung der Klasse new\_click\_tool\_app um die Verarbeitung von zwei Bildern zu ermöglichen
		\item new\_click\_tool\_main.cpp: Programmeinstieg. Initialisieren des Tools
		\item CMakeLists.txt: Build-Konfiguration des Subprojekts
	\end{itemize}
	
	\item \textbf{include/slam6d/fbr} enthält die Headerdatei new\_click\_tool\_app.h mit sämtlichen Deklarationen des Subprojekts
	\item In \textbf{bin} wird nach dem Build die ausführbare Datei abgelegt
\end{itemize}



\subsection{Genutzte Abhängigkeiten}

\begin{itemize}
    \item C++ Standardbibliothek
    \item stb\_image: Laden von Bilddateien
    \item Dear ImGui: Benutzungsoberfläche
    \item GLFW: Fensterverwaltung
    \item OpenGL: Darstellung von Bildern

\end{itemize}

\subsection{Einrichten der Entwicklungsumgebung}

Vorausgesetzt wird das Vorhandensein von
\begin{itemize}
	\item C++ 17-Compiler
	\item CMake (min. Version 3.14)
	\item 3DTK
\end{itemize}

\subsection{Build}
Das Tool ist vollständig im Hauptprojekt integriert und wird beim Build des Gesamtprojekts 3DTK miterstellt. Die Anleitung für den Build des Hauptprojekts befindet sich unter \url{https://github.com/JMUWRobotics/3DTK/blob/main/README.md}.
\vspace{1em}\\
Nach Erstellen der Makefiles durch \texttt{cmake} im Installationsprozess des Gesamtprojekts, kann das Tool auch unabhängig vom Gesamtprojekt kompiliert und genutzt werden. Um ausschließlich das Subprojekt \textit{new\_click\_tool} zu kompilieren, werden die Schritte der o.g. Installationsanweisung bis zum Aufruf von \texttt{make} ausgeführt. Anstelle von \texttt{make} wird \texttt{make new\_click\_tool} im Terminal ausgeführt.

\newpage

\section{Backend}


\subsection{Übersicht}

Das Backend des Tools \textit{new\_click\_tool} umfasst die Funktionalitäten zur Konvertierung von 3D-Dateiformaten in 2D-Bilder, die im Hauptfenster des GUIs angezeigt werden können. Weiterhin umfasst es Funktionen zur Speicherung sowie zum Laden gewählter Punkte bzw. Punktkorrespondenzen in Textdateien.

\subsection{Datei-/Ordnerstruktur}
Um Komplexität durch eine Vielzahl einzelner Dateien zu vermeiden, wurden die Backend-Funktionalitäten aufgrund ihres geringen Umfangs in der Hauptklasse \textit{new\_click\_tool\_app.cpp} und deren Header-Datei \textit{new\_click\_tool\_app.h} implementiert.


\subsection{Genutzte Abhängigkeiten}

\begin{itemize}
    \item 3DTK: Funktion \textit{scan\_to\_panorama} zur Konvertierung von 3D-Dateiformaten
    \item C++ Standardbibliothek
\end{itemize}

\subsection{Einrichten der Entwicklungsumgebung}
Vorausgesetzt wird das Vorhandensein von
\begin{itemize}
	\item C++ 17-Compiler
	\item CMake (min. Version 3.14)
	\item 3DTK
\end{itemize}

\subsection{Build}

Da die Funktionalitäten des Backends vollständig innerhalb der Frontend-Struktur implementiert wurden, werden diese gemeinsam mit dem Frontend kompiliert.

\newpage

\section{Bauen und Starten des gesamten Projekts}


\begin{itemize}
	\item 3DTK-Repository klonen
	\item Build des Gesamtprojekts gem. \href{https://github.com/JMUWRobotics/3DTK/blob/main/README.md}{Dokumentation}
	\item Speicherort des Tools: \textbf{3DTK/bin/}
	\item Start des Projekts über Binary-Datei oder
	\item Kommandozeile \hspace{2em} z.B.:\\
\begin{tabular}{p{9.3cm} p{5.7cm}}
		\texttt{./new\_click\_tool} & \textit{Ohne geladene Bilder}\\
		
		\texttt{./new\_click\_tool -help} & \textit{Hilfestellung zum Aufruf des Tools}\\
		
		\texttt{./new\_click\_tool -im <Bildpfad>} & \textit{Mit einem geladenen Bild}\\
		
		\texttt{./new\_click\_tool -scan <Scanpfad>} & \textit{Mit einem geladenen Scan in default format (uos)}\\

		\texttt{./new\_click\_tool -scan <Scanpfad> -f <Dateiformat>} & \textit{Mit einem geladenen Scan mit einem 3DTK-spezifischen Format}\\
		
		\texttt{./new\_click\_tool -im <Bildpfad> -im <Bildpfad>} & \textit{Mit zwei geladenen Bildern}\\
		
		
		\texttt{./new\_click\_tool -im <Bildpfad> -out <Speicherort>} & \textit{Mit einem geladenen Bild und manuell definiertem Speicherort für erzeugte Panoramabilder und Koordinaten}\\
		
	\end{tabular}
\end{itemize}


\end{document}